% Build here: latexmk -pdf -outdir=../../../build/audits/rfp_intrinsic a4_associator.tex
\documentclass{article}
\usepackage[a4paper,margin=1in]{geometry}
\usepackage{amsmath,amssymb,amsthm,mathtools}
\usepackage{booktabs}
\newtheorem{theorem}{Theorem}
\newtheorem{lemma}[theorem]{Lemma}
\newtheorem{remark}[theorem]{Remark}
\newcommand{\C}{\mathbb C}
\newcommand{\Hom}{\operatorname{Hom}}
\newcommand{\id}{\mathrm{id}}
\newcommand{\diag}{\operatorname{diag}}
\newcommand{\Span}{\operatorname{span}}
\begin{document}

\begin{center}
{\Large The exact weighted associator test for $\operatorname{Rep}(A_4)$}\\
{\large Classification of the two triplet-channel weights}
\end{center}

\paragraph{Status.}
\textbf{SOLVED.}  For the stated diagonal path weight
\[
\Lambda(p,q)=\diag(1,1,1,p^2,pq,pq,q^2),
\]
the symmetric--antisymmetric Clebsch--Gordan basis gives
$[F,\Lambda(p,q)]=0$ exactly when $p=q=1$ or $p=q=-1$.
After allowing an arbitrary unitary basis change in the unique multiplicity
space $\Hom_{A_4}(3,3\otimes3)\cong\C^2$, there is one additional unequal
branch:
\[
\{p,q\}=\{i,-i\}.
\]
It is realized by replacing the symmetric and antisymmetric channels with the
two oriented channels. There are no other complex solutions. Object
rescalings, one-dimensional Clebsch--Gordan phases, and nonzero BNT basis
rescalings do not enlarge this list because the associativity equation is
covariant under all such changes.

\section{Precise formulation}

Let
\[
A_4=\langle s,t\mid s^2=t^3=(st)^3=1\rangle.
\]
Let $V=\C^3$ with orthonormal basis $e_0,e_1,e_2$, indices understood modulo
$3$, and triplet representation
\[
\rho(s)=\diag(1,-1,-1),
\qquad
\rho(t)e_j=e_{j+1}.
\]
Put $\omega=e^{2\pi i/3}$. The one-dimensional representations are
$\chi_r(s)=1$ and $\chi_r(t)=\omega^r$ for $r=0,1,2$; they are denoted
$1,1',1''$, respectively. The tensor product rule is
\[
V\otimes V\cong \chi_0\oplus\chi_1\oplus\chi_2\oplus V_s\oplus V_a.
\]

The seven-dimensional path space for $V,V,V\to V$ has ordered basis
\[
(1,1',1'',ss,sa,as,aa),
\]
where, for example, $sa$ means that the inner pair uses the symmetric triplet
embedding and the outer fusion uses the antisymmetric triplet embedding.
The same order is used for left- and right-associated paths. If $W_L$ and
$W_R$ are the corresponding orthonormal families of intertwiners
$V\to V^{\otimes3}$, the recoupling matrix is
\[
F=W_R^*W_L,
\]
where the adjoint uses the normalized Hilbert--Schmidt inner product
$\langle X,Y\rangle=\frac13\operatorname{tr}(X^*Y)$ on the intertwiner space.

The candidate weight operator on the multiplicity space
$M=\Hom_{A_4}(V,V\otimes V)$ has eigenvalues $p,q$. In an eigenbasis its
induced operator on the four triplet paths is
$\diag(p,q)\otimes\diag(p,q)$. The three one-dimensional intermediate paths
have weight one. Thus the associator test in that basis is
\[
F\Lambda(p,q)=\Lambda(p,q)F.
\]
A solution means that this equality holds after a coherent choice of basis of
$M$. Coherent means that the same basis of the same fusion space $M$ is used
at every occurrence of a $3,3\to3$ vertex.

\section{Unitary Clebsch--Gordan embeddings}

For $r=0,1,2$, define the isometric embedding
$E_r:\chi_r\to V\otimes V$ by
\[
E_r(1)=\frac1{\sqrt3}\sum_{j=0}^2\omega^{-rj}e_j\otimes e_j.
\]
Define the two isometric triplet embeddings $C_s,C_a:V\to V\otimes V$ by
\begin{align*}
C_s e_j&=\frac1{\sqrt2}
 (e_{j+1}\otimes e_{j+2}+e_{j+2}\otimes e_{j+1}),\\
C_a e_j&=\frac1{\sqrt2}
 (e_{j+1}\otimes e_{j+2}-e_{j+2}\otimes e_{j+1}).
\end{align*}
Finally, identify $\chi_r\otimes V$ and $V\otimes\chi_r$ with copies of $V$
and use
\[
J_r e_j=\omega^{rj}(1\otimes e_j),
\qquad
\widetilde J_r e_j=\omega^{rj}(e_j\otimes1).
\]
Each displayed map is an $A_4$-intertwining isometry. The generator checks
are immediate from the cyclic action of $t$ and the signs of $s$; explicitly,
\[
(\rho\otimes\rho)(g)E_r=E_r\chi_r(g),
\qquad
(\rho\otimes\rho)(g)C_x=C_x\rho(g),
\]
for $x=s,a$ and $g=s,t$, with analogous identities for $J_r$ and
$\widetilde J_r$.

Define the left paths
\begin{align*}
L_r&=(E_r\otimes\id_V)J_r, &&r=0,1,2,\\
L_{xy}&=(C_x\otimes\id_V)C_y, &&x,y\in\{s,a\},
\end{align*}
and the right paths
\begin{align*}
R_r&=(\id_V\otimes E_r)\widetilde J_r, &&r=0,1,2,\\
R_{xy}&=(\id_V\otimes C_x)C_y, &&x,y\in\{s,a\}.
\end{align*}
Direct contraction gives
\[
L_\alpha^*L_\beta=R_\alpha^*R_\beta
 =\delta_{\alpha\beta}\id_V.
\]
Thus both families are orthonormal bases of
$\Hom_{A_4}(V,V^{\otimes3})$, whose dimension is
$3+2\cdot2=7$.

\section{The exact recoupling matrix}

In the path order $(1,1',1'',ss,sa,as,aa)$, direct contraction gives
\[
F=
\begin{pmatrix}
 \frac13&\frac13&\frac13&\frac1{\sqrt3}&0&0&-\frac1{\sqrt3}\\
 \frac13&\frac13&\frac13&-\frac1{2\sqrt3}&-\frac i2&\frac i2&\frac1{2\sqrt3}\\
 \frac13&\frac13&\frac13&-\frac1{2\sqrt3}&\frac i2&-\frac i2&\frac1{2\sqrt3}\\
 \frac1{\sqrt3}&-\frac1{2\sqrt3}&-\frac1{2\sqrt3}&\frac12&0&0&\frac12\\
 0&-\frac i2&\frac i2&0&-\frac12&-\frac12&0\\
 0&\frac i2&-\frac i2&0&-\frac12&-\frac12&0\\
 -\frac1{\sqrt3}&\frac1{2\sqrt3}&\frac1{2\sqrt3}&\frac12&0&0&\frac12
\end{pmatrix}.
\]
It satisfies
\[
F^*F=I_7,
\qquad
W_L=W_R(F\otimes I_V).
\]

For the unrotated symmetric--antisymmetric basis, the nonzero entries of
$F\Lambda-\Lambda F$ are scalar multiples of
\[
p^2-1,\qquad pq-1,\qquad q^2-1,\qquad p^2-q^2.
\]
For example,
\begin{align*}
(F\Lambda-\Lambda F)_{1,4}&=\frac{p^2-1}{\sqrt3},\\
(F\Lambda-\Lambda F)_{2,5}&=-\frac i2(pq-1),\\
(F\Lambda-\Lambda F)_{1,7}&=-\frac{q^2-1}{\sqrt3}.
\end{align*}
Consequently this particular basis gives
\[
[F,\Lambda(p,q)]=0
\quad\Longleftrightarrow\quad
p^2=pq=q^2=1
\quad\Longleftrightarrow\quad
p=q\in\{1,-1\}.
\]
This calculation alone is not basis-independent, because a non-diagonal
unitary change of the multiplicity basis also rotates the local weight
operator. The next section performs the invariant test.

\section{Classification under every coherent $U(2)$ change}

Write the recoupling matrix in $3+4$ block form
\[
F=\begin{pmatrix}A&B\\B^*&D\end{pmatrix}.
\]
In the ordered triplet-path basis $(ss,sa,as,aa)$, introduce
\begin{align*}
r_1&=(1,0,0,-1)^T,&
r_2&=(0,1,-1,0)^T,\\
n_1&=(0,1,1,0)^T,&
n_2&=(1,0,0,1)^T.
\end{align*}
The displayed matrix gives
\[
\operatorname{ran}(B^*)=\Span\{r_1,r_2\},
\qquad
\ker B=\Span\{n_1,n_2\},
\]
and
\[
Dr_1=Dr_2=0,
\qquad
Dn_1=-n_1,
\qquad
Dn_2=n_2.
\]

Let $U\in U(2)$ be an arbitrary new basis of the multiplicity space and put
\[
H=U\diag(p,q)U^*,
\qquad
K=H\otimes H.
\]
The transformed associator commutes with the diagonal path weight if and only
if
\[
\left[F,I_3\oplus K\right]=0.
\]
The off-diagonal blocks imply
\[
Kr_1=r_1,
\qquad
Kr_2=r_2,
\]
and the lower-right block implies $[D,K]=0$.

\begin{lemma}
Let $H\in M_2(\C)$. Then
\[
[ F,I_3\oplus(H\otimes H)]=0
\]
if and only if
\[
H\in\{I_2,-I_2,i\sigma_x,-i\sigma_x\},
\qquad
\sigma_x=\begin{pmatrix}0&1\\1&0\end{pmatrix}.
\]
\end{lemma}

\begin{proof}
Write
\[
H=\begin{pmatrix}a&b\\c&d\end{pmatrix}.
\]
The equations $(H\otimes H)r_1=r_1$ and
$(H\otimes H)r_2=r_2$ give
\begin{align}
a^2-b^2&=1, & ac-bd&=0, & c^2-d^2&=-1,
\label{eq:r1}\\
ad-bc&=1.\label{eq:r2}
\end{align}
Selected entries of $[D,H\otimes H]=0$ give
\begin{align*}
3ab+cd&=0,&ab+3cd&=0,\\
3ac+bd&=0,&ac+3bd&=0.
\end{align*}
Hence
\[
ab=cd=ac=bd=0.
\]
If $a\ne0$, then $b=c=0$. Equations
\eqref{eq:r1}--\eqref{eq:r2} give $a=d=1$ or $a=d=-1$, so
$H=\pm I_2$. If $a=0$, then \eqref{eq:r2} gives $bc=-1$, so
$b,c\ne0$; the equations $bd=cd=0$ imply $d=0$. Equation
\eqref{eq:r1} then gives $b^2=c^2=-1$ and $bc=-1$, hence
$b=c=i$ or $b=c=-i$. Thus $H=\pm i\sigma_x$.

Conversely, direct substitution shows that all four matrices satisfy the
commutator equation.
\end{proof}

\begin{theorem}[Exact classification]
Let $p,q\in\C$. There exists a coherent unitary basis of
$\Hom_{A_4}(3,3\otimes3)$ for which the weighted associator commutes with
$\Lambda(p,q)$ if and only if
\[
(p,q)=(1,1),\quad(-1,-1),\quad(i,-i),\quad(-i,i).
\]
Equivalently, the unordered spectrum is
\[
\{p,q\}=\{1,1\},\quad\{-1,-1\},\quad\{i,-i\}.
\]
In particular, the unique unequal branch is
\[
p=-q,\qquad p^2=-1.
\]
\end{theorem}

\begin{proof}
The matrix $H=U\diag(p,q)U^*$ has spectrum $\{p,q\}$. Apply the lemma.
The spectra of $I_2,-I_2,i\sigma_x,-i\sigma_x$ are, respectively,
$\{1,1\}$, $\{-1,-1\}$, and $\{i,-i\}$ for the last two cases.
\end{proof}

\section{The exceptional oriented-channel solution}

Take
\[
U=\frac1{\sqrt2}\begin{pmatrix}1&1\\1&-1\end{pmatrix},
\qquad
\diag(p,q)=\diag(i,-i).
\]
Then
\[
U\diag(i,-i)U^*=i\sigma_x.
\]
The new Clebsch--Gordan embeddings are
\begin{align*}
C_+&=\frac{C_s+C_a}{\sqrt2},&
C_+e_j&=e_{j+1}\otimes e_{j+2},\\
C_-&=\frac{C_s-C_a}{\sqrt2},&
C_-e_j&=e_{j+2}\otimes e_{j+1}.
\end{align*}
Thus the exceptional basis consists of the two orientations of the ordered
pair complementary to $j$. In the induced path basis, the recoupling matrix
is
\[
F_U=(I_3\oplus U\otimes U)^*F(I_3\oplus U\otimes U)
\]
with
\[
F_U=
\begin{pmatrix}
\frac13&\frac13&\frac13&0&\frac1{\sqrt3}&\frac1{\sqrt3}&0\\
\frac13&\frac13&\frac13&0&-\frac1{2\sqrt3}+\frac i2&-\frac1{2\sqrt3}-\frac i2&0\\
\frac13&\frac13&\frac13&0&-\frac1{2\sqrt3}-\frac i2&-\frac1{2\sqrt3}+\frac i2&0\\
0&0&0&0&0&0&1\\
\frac1{\sqrt3}&-\frac1{2\sqrt3}+\frac i2&-\frac1{2\sqrt3}-\frac i2&0&0&0&0\\
\frac1{\sqrt3}&-\frac1{2\sqrt3}-\frac i2&-\frac1{2\sqrt3}+\frac i2&0&0&0&0\\
0&0&0&1&0&0&0
\end{pmatrix}.
\]
For $(p,q)=(i,-i)$,
\[
\Lambda(i,-i)=\diag(1,1,1,-1,1,1,-1),
\]
and the zero pattern of $F_U$ makes
$[F_U,\Lambda(i,-i)]=0$ immediate.

\section{Gauge and rescaling checks}

A general change of left and right path coordinates has the form
\[
W_L\mapsto W_LG_L,
\qquad
W_R\mapsto W_RG_R,
\]
with invertible $G_L,G_R$. The recoupling matrix and the two path-weight
operators transform as
\[
F\mapsto G_R^{-1}FG_L,
\qquad
\Lambda_L\mapsto G_L^{-1}\Lambda_LG_L,
\qquad
\Lambda_R\mapsto G_R^{-1}\Lambda_RG_R.
\]
Therefore
\[
F\Lambda_L=\Lambda_RF
\]
is equivalent to the transformed equation. This covers unitary
multiplicity-basis changes, phases in one-dimensional fusion spaces, and
nonzero rescalings of BNT representatives. Such coordinate changes cannot
turn a failed associativity equation into a valid one.

If each simple object $a$ is rescaled by a nonzero scalar $z_a$, a binary
coefficient for $a\otimes b\to c$ acquires the coboundary factor
$z_az_b/z_c$. Along either association of $a\otimes b\otimes c\to d$, the
product of the two factors is
\[
\frac{z_az_b}{z_x}\frac{z_xz_c}{z_d}
 =\frac{z_az_bz_c}{z_d},
\]
independent of the intermediate object $x$. Hence every diagonal entry of
the path-weight operator receives the same overall scalar, which cancels
from the commutator test. Object rescaling does not alter the classification.

\begin{remark}[Occurrence-dependent bases are not a fusion gauge]
The unique multiplicity space is $M=\Hom_{A_4}(3,3\otimes3)$. A fusion datum
chooses one basis of $M$, used at every trivalent vertex of that type. If one
instead assigns unrelated $U(2)$ bases to separate occurrences of the same
vertex while keeping the numerical diagonal matrix $\diag(p,q)$ fixed at
each occurrence, one has changed the local weight operator from occurrence
to occurrence. That is a different, nonlocal datum rather than a gauge
transformation of a single fusion rule. The classification above uses the
standard coherent convention.
\end{remark}

\section{Verification record}

The accompanying deterministic script
\texttt{docs/audits/rfp\_intrinsic/checks/a4\_associator.py} performs the Clebsch--Gordan
contractions over the exact field $\mathbb Q(i,\sqrt3)$. It verifies:
\begin{enumerate}
\item the relations $s^2=t^3=(st)^3=I$;
\item equivariance of all fourteen left and right path embeddings under both
generators;
\item exact path orthonormality $W_\alpha^*W_\beta=\delta_{\alpha\beta}I_3$;
\item the identity $W_L=W_R(F\otimes I_3)$;
\item unitarity $F^*F=I_7$;
\item every entry of the fixed-basis symbolic commutator.
\end{enumerate}
The arbitrary-$H$ classification and the exceptional Hadamard-gauge check
used in this report were recomputed independently of the script;
they are derived in the preceding sections but are not implemented in the
checked-in script. The script writes no files.

\paragraph{Formalization target.}
The most useful self-contained lemma for Lean 4 is the $2\times2$ matrix
classification
\[
[ F,I_3\oplus(H\otimes H)]=0
\Longleftrightarrow
H\in\{\pm I_2,\pm i\sigma_x\}.
\]
It only requires finite matrix extensionality and elementary polynomial
algebra; it is delicate enough to merit formalization because it is the step
that excludes all hidden multiplicity gauges.

\end{document}